\section{Method}
\label{sec:method}
\subsection{Problem Formulation}
Let $G = (V, E, \phi, \psi, \tau)$ be a heterogeneous temporal graph, where
$\phi: V \to \mathcal{A}$ maps nodes to a set of types $\mathcal{A}$,
$\psi: E \to \mathcal{R}$ maps edges to relation types $\mathcal{R}$, and
$\tau: E \to \mathbb{R}$ assigns each edge a timestamp. Given a query node
$v$ at reference time $t$, \sysname{} produces an embedding
$z_v^t \in \mathbb{R}^d$ by aggregating over the time-decayed neighborhood
$N_t(v) = \{u : (u,v) \in E,\ \tau(u,v) \le t\}$.

\subsection{Type-Aware Encoder}
Each node type $a \in \mathcal{A}$ owns a projection matrix
$W_a \in \mathbb{R}^{d \times d_{\text{in}}}$. A neighbor $u$ of type
$\phi(u)$ contributes a message $m_u = W_{\phi(u)} x_u$, where $x_u$ is its
input feature vector. Messages are combined with a shared multi-head
attention layer so that type-specific projections stay disjoint while the
attention weights are learned jointly across types, letting the model
transfer attention patterns learned on data-rich types (e.g.\ \emph{item})
to data-sparse ones (e.g.\ \emph{seller}).

\subsection{Time-Decayed Sampling}
Neighbors are drawn with sampling probability
$p(u \mid v, t) \propto \exp\!\big(-\lambda\,(t - \tau(u,v))\big)$,
where $\lambda$ controls how quickly influence decays; we set $\lambda$
per-dataset via the validation split (Section~\ref{sec:experiments}). This
keeps the neighborhood size bounded — Figure~\ref{fig:arch} sketches the
resulting pipeline from raw edges to the final embedding.

\begin{figure}[t]
\centering
\begin{tikzpicture}[
  node distance=6mm and 6mm,
  box/.style={draw=kddblue,fill=kddblue!8,rounded corners=2pt,
              minimum height=8mm,minimum width=27mm,align=center,
              font=\footnotesize},
  arr/.style={-{Stealth[length=2mm]},kddblue,thick}
]
\node[box] (input) {Heterogeneous\\temporal edges};
\node[box, below=of input] (sample) {Time-decayed\\neighbor sampler};
\node[box, below=of sample] (encode) {Type-aware\\attention encoder};
\node[box, below=of encode, fill=kddgold!15, draw=kddgold] (embed) {Embedding $z_v^t$};
\draw[arr] (input) -- (sample);
\draw[arr] (sample) -- (encode);
\draw[arr] (encode) -- (embed);
\end{tikzpicture}
\caption{\sysname{} pipeline: edges are sampled with exponential time decay,
routed through type-specific projections, then fused by shared attention.}
\label{fig:arch}
\end{figure}

\subsection{Contrastive Objective}
For an anchor interaction between node $i$ and a temporally adjacent
partner $i^+$, and a batch of $N$ negatives drawn uniformly from the same
type, we minimize the InfoNCE-style loss
\begin{equation}
\mathcal{L}_{\text{con}} = -\log
\frac{\exp\!\big(\mathrm{sim}(z_i, z_{i^+}) / \tau_s\big)}
     {\sum_{j=1}^{N} \exp\!\big(\mathrm{sim}(z_i, z_j) / \tau_s\big)},
\label{eq:contrastive}
\end{equation}
where $\mathrm{sim}(\cdot,\cdot)$ is cosine similarity and $\tau_s$ is a
learned temperature, initialized at $\tau_s = 0.07$. Sampling positives from
real temporal adjacency, rather than synthetic augmentation, means the
objective directly reflects the recency structure the downstream tasks care
about. Encoding a node's neighborhood costs $O(k \log k)$ for a sampled
neighborhood of size $k$, dominated by the decay-weighted sort; end-to-end
training on the largest benchmark graph runs in $O(|E| \log k)$.
